%% Section 7 -- Discussion & Operational Considerations
%% Ported from docs/DRAFT_PAPER.md Section 7.
%%
%% Section 7.4's first limitation is the Section 8.2 contingency
%% (docs/DRAFT_PAPER.md drafting note 3). Its wording is deliberate and was NOT
%% softened during LaTeX assembly. Do not soften it.

\section{Discussion \& Operational Considerations}
\label{sec:discussion}

\subsection{End-to-end encryption}
\label{sec:discussion:encryption}

CARTA operates on decrypted local exports that the user's own device produced,
using the platform's own export function, on an account the user controls.
Encryption is not an obstacle the pipeline encounters, because the pipeline begins
after the point where the platform has already decrypted the content for its
owner. Nothing in Sections~\ref{sec:carta} or~\ref{sec:algorithms} attempts to
read encrypted backup blobs, recover keys, or access an account other than the one
performing the export. This is worth stating plainly because reconstruction of chat history is
adjacent to a body of forensic work that does address encrypted stores, and a
reader arriving from that direction may reasonably expect such a claim. We make
none.

The corollary is that CARTA inherits whatever the export function chose to
include. A platform that omits a field from its export puts that field beyond
reach of any conforming implementation, and reconstruction cannot recover
information the export never carried.

\subsection{Generalisation across platforms}
\label{sec:discussion:generalisation}

The contract \carta{} is stated without reference to any platform, and
Normalization, Synthesis, the invariants and the report structure are defined over
record types rather than over file formats. Ingestion is the exception, and it is
necessarily platform-specific: Signal, Telegram and iMessage each export in their
own form, with their own timestamp conventions, their own representation of system
events, and their own attachment references.

Whether the rest of the contract survives contact with those formats is an open
question, and we do not claim it has been answered. The specification has been
developed against one platform's export conventions, illustrated by a corpus drawn
from that same platform. A second platform could plausibly require reason codes
the taxonomy in Section~\ref{sec:carta:normalization} does not anticipate, or
exhibit a form of cross-fragment disagreement that the conflict definition in
Section~\ref{sec:carta:synthesis} does not capture. Establishing generalisation
requires implementing Ingestion for at least one further platform and checking
whether the downstream stages need amendment. That work has not been done.

\subsection{Forensic and legal considerations}
\label{sec:discussion:forensic}

CARTA's provenance model and its reconstruction report are close in shape to
chain-of-custody practice, which Section~\ref{sec:related:forensics} discusses:
each record traces to an identified source with a byte range and a content hash,
originals are read and never modified, and the procedure emits an account of what
it did. The resemblance is architectural and deliberate.

It is not a claim of admissibility. This paper offers no legal analysis, no
jurisdiction-specific argument, and no forensic validation of any kind, and none
of the three would follow from the specification alone. Establishing that a
reconstruction produced under CARTA would be accepted as evidence anywhere would
require an implementation, validation against known-ground-truth material, and
legal argument that is outside both this paper's scope and its author's competence
to assert. The honest position is that the design borrows a discipline from
forensic practice without claiming the standing that practice confers.

%% Added 2026-08-26. Section 2.2 promises that this section treats fragments
%% which are not internally complete; before this subsection existed, that was a
%% dangling forward reference (peer review, Minor 5). Do not delete without also
%% removing the promise in Section 2.2's internal-completeness paragraph.
\subsection{Fragments that are not internally complete}
\label{sec:discussion:incomplete}

Section~\ref{sec:tcf:conditions} defines internal completeness as a fragment
being well-formed and gap-free with respect to its own accumulation window, and
says plainly that this is an idealisation of the TCF model rather than a property
established of any real archive. Real archives violate it. A user clears a
conversation, a device runs out of storage mid-export, a chat is deleted on one
device and not the other, or an export is taken from an application that had
already pruned old media. The question this raises is what the specification owes
a reader when its own precondition fails.

The answer follows from what CARTA does and does not attempt. Nothing in
Sections~\ref{sec:carta} or~\ref{sec:algorithms} tests whether a fragment is
internally complete, and nothing in the pipeline depends on its being so. The
stages consume records; a fragment with a hole in it simply contributes fewer
records, and every stage behaves exactly as specified on the ones it does
contribute. $\mathcal{U}$ is still the deduplicated union of what arrived,
$\Omega$ still accounts for every record that did not reach it, and Invariants 1
and 2 are unaffected, because both are claims about records the pipeline saw
rather than about records that existed once.

What is lost is an interpretive guarantee, not a mechanical one. Where both
fragments are internally complete, a gap in $\mathcal{U}$ over some interval
licenses the inference that no messages were sent in it. Where they are not, that
inference is unavailable, and $\mathcal{U}$ cannot distinguish a quiet week from a
cleared conversation. This is the sharpest limitation in the paper, because it is
invisible in the output: nothing in $\mathcal{U}$ or $\Omega$ marks the
difference, and a reader would have to already know the fragment's history to
suspect it.

Closing this properly would mean detecting absence, which is a genuinely harder
problem than reconciling presence and is not attempted here. Export formats do
not record what they omitted. A partial approach available to an implementation
is to report per-fragment coverage in $\Omega_{\text{stats}}$ (first and last
record timestamps per fragment, and the distribution of inter-record intervals),
which would let a reader see that one fragment claims a window it contributes
almost nothing to. That is a heuristic surfaced for inspection rather than a
completeness test, which is the most this specification's discipline permits it to
be, and stating it as future work is more honest than implying the model's
idealisation is harmless.

%% Added 2026-08-29. This subsection is the main scientific yield of building the
%% implementation, and each item is a place the SPECIFICATION was wrong, not a
%% place the code was. Do not soften these into "clarifications": the point is
%% that a document two review rounds had gone over still contained them.
\subsection{What implementing the specification revealed}
\label{sec:discussion:specfindings}

Writing a conforming implementation surfaced four defects in the specification
itself. Three of them survived two rounds of review of the prose. They are
recorded here because a specification paper that reported none would be making a
weaker claim than it appears to: the value of building the reference
implementation was mostly this, and not the reconstruction counts.

\paragraph{A UTC timestamp with no offset to convert by.}
Section~\ref{sec:carta:normalization} requires a canonical record to carry a UTC
epoch timestamp. A WhatsApp text export carries no UTC offset.
Section~\ref{sec:implementation:inputs} already observes that the export has no
locale marker to settle the date ordering, and the same absence applies to the
clock, so converting a wall-clock reading into an instant needs information the
input does not contain. Resolving it silently is the failure Normalization exists
to prevent. The implementation makes the offset a declared parameter defaulting
to zero, so canonical timestamps are wall-clock instants correct up to a stated
offset, and the statement reaches the reader through
$\Omega$. Reconciling fragments from devices in different zones requires setting
it per fragment. The specification should say where the offset comes from; as
written it asserts UTC and leaves the question unasked.

\paragraph{Equation~(\ref{eq:conflict}) reports ordinary conversation as
contradiction.} Its second disjunct is $s(u_i) \neq s(u_j)$. Read literally, any
two records in the same coarse bucket sent by different people conflict. Three
participants messaging inside one minute produces three conflict entries, and the
count grows with the square of the traffic in a bucket, so
$\Omega_{\text{conflicts}}$ would be dominated by ordinary group-chat activity on
any real archive. Two people talking is not an irreconcilable observation.

The reading the implementation declares instead is that a sender disagreement is
irreconcilable only when the two records are otherwise the same observation:
same coarse bucket, not deduplicated, same body, different sender. That is a
genuine question of who said something, and no total order settles it. Records
with no sender are excluded from the clause, because $s(\cdot)$ is undefined for
a system line and comparing it against a name produces a finding about a type
error rather than about the archive. The ordering clause is unaffected.

\paragraph{One export variant defeats deduplication.}
Equation~(\ref{eq:dedup}) matches on body equality. The doubled-filename
attachment variant of Section~\ref{sec:implementation:inputs} repeats the
filename on a continuation line, which becomes part of the body. Two devices
exporting the same attachment message in different variants therefore produce
bodies that differ by that line, fail to match, and place the message in
$\mathcal{U}$ twice. This is not a corner case for TCF specifically: two
handsets on possibly different application versions is the configuration the
whole paper is about. The implementation canonicalises the repeated line away in
Normalization, where converting an export artefact into a canonical payload
belongs, which leaves Equation~(\ref{eq:dedup}) intact rather than widening it to
tolerate export noise.

\paragraph{An unreachable branch, and what it says about the reason codes.}
The fragment-scope step was written to re-quarantine a record when the
established ordering yields an impossible date. That branch cannot execute. A
record is quarantined as ambiguous only when both leading fields lie in 1 to 12,
which is a valid day-and-month pair under either reading, and its clock has
already been validated before the ordering question arises. The case is a guard
rather than a behaviour. It is worth recording because it confirms the
distinction Section~\ref{sec:carta:normalization} draws between
\texttt{MALFORMED\_\allowbreak TIMESTAMP} and
\texttt{AMBIGUOUS\_\allowbreak DATE\_\allowbreak FORMAT} is exhaustive on this
input format: a date string is unreadable, or readable one way, or readable two
ways, and the third case never collapses into the first.

\subsection{Limitations}
\label{sec:discussion:limitations}

\paragraph{The related-work position is contingent on two documents we could not
obtain.}
Section~\ref{sec:related:forensics} argues that the forensic lineage stages
reconstruction without specifying a typed inter-stage contract. That argument
rests on the sources we were able to read. Two in the same lineage resisted
retrieval, and either could narrow the claim further, in which case the
contribution reduces to the quarantine set and the report as the properties
without precedent.

\paragraph{The clustering procedure is deliberately unspecified.}
Section~\ref{sec:algorithms:similarity} fixes the similarity criterion and leaves
the grouping algorithm open, and transitive closure with a bounded diameter and
connected components over a similarity graph can disagree on chained
near-boundary records. Two conforming implementations may therefore produce
different reconciliation clusters from the same input while each remaining
internally deterministic. This is a real gap in the specification's strength, and
closing it means choosing a rule and defending the choice, which we have not done.

\paragraph{Determinism and idempotence are requirements, not demonstrated
properties.}
Invariant 3 constrains an implementation to reproduce its own output, and
Section~\ref{sec:evaluation:reconstruction} reports one implementation doing so
on one input across two consecutive runs. That is an instance, not a proof, and
nothing establishes it holds for inputs the implementation has not seen.
Idempotence is weaker
still, which is why Section~\ref{sec:carta:invariants} declines to number it among
the invariants: it is stated as prose about what Synthesis must not do, and the
type-safe formalisation establishing that records in $\mathcal{U}$ are
schema-compatible with those in $\mathcal{R}'$ is deferred rather than supplied.

\paragraph{The corpus is illustrative, synthetic and single-platform.}
Nineteen edge-case categories across five synthetic files show that the input
space is varied. They do not establish coverage of that space, no claim of
completeness attaches to the corpus, and, as
Section~\ref{sec:implementation:inputs} states, the categories record what this
specification anticipates rather than what WhatsApp has been observed to emit.

\paragraph{The platform's reunification affordances were not systematically
enumerated.}
Condition 3 asserts that no platform-native operation yields the union of two
archives, and the support offered for it is documentary: the vendor's restore
guidance describes no such operation, and third-party merge tools exist to supply
what it describes no way of doing. We did not enumerate the platform's export,
restore, transfer and migration paths and record what each does to a pair of
archives. That enumeration is available to anyone holding the configuration, and
it would convert the weakest premise in this paper into its only piece of primary
evidence. It would be a dated, versioned, single-case observation rather than an
experiment, and worth reporting as exactly that. We report its absence rather than implying the documentary
argument is stronger than it is.
